\section{Federated Learning Algorithms}
In this section, we delve into the foundational algorithmic implementation that has shaped Federated Learning
research, by introducing a structural background for subsequent efforts. This, algorithm is none other than
\emph{FedAvg}, which introduced several key innovations, that contributed as the basis for more sophisticated
optimized approaches that are currently employed in the domain. Nonetheless, Before exploring the novelties of
FedAvg, we will first discuss FedSGD, as the simplest Federated Learning Algorithm that set the stage for
the FedAvg deployement. Concluding the analysis of the two pricncipal algorrithms, the shortcomings of contemporary
FL implementations are highligthed.

Before focusing on the algorithmic implementations a common notation ground is laid with the \cref{tab:fl_notation}:
\begin{table}[h!]
    \centering
    \begin{tabular}{|c|p{10cm}|}
        \hline
        \rowcolor{lightgray}
        \textbf{Notation} & \textbf{Definition}                                                                \\
        \hline
        $w_t$             & Model weights on communication round \# $t$                                        \\
        % \hline
        $w_t^k$           & Model weights on communication round \# $t$ on client $k$                          \\
        % \hline
        $K$               & Number of all clients                                                              \\

        $C$               & Fraction of clients performing computations in each round                          \\
        % \hline
        $E$               & Number of training passes each client makes over its local dataset on each round   \\
        % \hline
        $B$               & The local minibatch size used for the client updates                               \\
        % \hline
        $\eta$            & The learning rate                                                                  \\
        % \hline
        $\mathcal{P}_k$   & Set of data points on client $k$                                                   \\
        % \hline
        $n_k$             & Number of data points on client $k$                                                \\
        % \hline
        $f_i(w)$          & Loss $l(x_i, y_i; w)$ i.e., loss on example $(x_i, y_i)$ with model parameters $w$ \\
        \hline
    \end{tabular}
    \caption{Notation used in Federated Learning algorithms}
    \label{tab:fl_notation}
\end{table}

%%%%%%%%%%%%%%%%%%%%%%%%%%%%%%%%%%%%%%%%%%%%%%%%%%%%%%%%%%%%%%%%%%%%%%%%%%%%%%%%%%%%%%%%%%%%%%%%%%%%%%%%%%%%%%%%%%%%%%%%%%%%%%%%%%%%%%%%%%%%%%%%%%%%%%%%%%%%%%
\subsection{Federated Stochastic Gradient Descent (FedSGD)}
Stochastic Gradient Descent (SGD), has shown great results in deep learning, making it a foundational algorithm for many
optimization problems. Given its effectiveness, researchers decided to base the Federated Learning training algorithm on
SGD as a baseline. In traditional Stochastic Gradient Descent,
% Mini-Batch Gradient Descent, or for simplicity SGD,
% \footnote{As stated ealier Mini-Batch Gradient Descent and Stochastic Gradient Descent are commonly used interchangeably},
gradients are computed on a random subset of the total dataset and then used to
update the model parameters. This process is repeated iteratively, making it computationally efficient for centralized learning.

However, applying SGD naively to the federated optimization problem involves a single batch gradient calculation on a randomly
selected client per round of communication. This way each selected client participates in training using the local data points.
While this approach is computationally efficient, it requires a large number of
communication rounds to achieve good model performance. This is due to the fact that each client only contributes a small fraction
of the data in each round.

Federated Stochastic Gradient Descent is a direct adaptation of the traditional SGD algorithm
\footnote{Or if we want to be precise, it is a direct adaptation of Mini-Batch Gradient Descent}
to the federated setting. In the FedSGD process, the server starts by initializing the global model weights \( w_t \). For each
communication round \( t \), a fraction \( C \) of clients is randomly selected to participate. Each selected client \( k \) uses its
local dataset \( \mathcal{P}_k \) to compute the gradient of the loss function \( f_i(w) \) over all its data points \( n_k \). These
gradients \( G_k \) are then sent to the central server, where they are aggregated. The server updates the global model weights \( w_t \)
by taking a weighted average of the received gradients, with weights proportional to the number of data points \( n_k \) on each client.
The updated global model \( w_{t+1} \) is then distributed back to all clients, and the process repeats until the stopping criterion is met.

\begin{algorithm}[H]
    \caption{Federated Stochastic Gradient Descent (FedSGD)}
    \label{alg:fedsgd}
    \begin{algorithmic}[1]
        \State \textbf{Server Execution:}
        \State Initialize global model weights $w_0$, and learning rate $\eta$
        \For{each round $t=1, 2, \ldots$}
        \State $m \gets \max(C \cdot K, 1)$, number of participating clients
        \State $S_t \gets$ random set of $m$ clients
        \State Send global model $w_t$ to $S_t$ clients
        \For{each client $k \in S_t$ \textbf{in parallel}}
        \State $g_k \gets$ \textsc{GradientStep}$(k, w_t)$
        \EndFor
        \State $w_{t+1} \gets w_t - \eta \sum_{k \in S_t} \frac{n_k}{n} g_k$
        \State Send the model $w_{t+1}$ to all participants
        \EndFor

        \vspace{1em}
        \State \textbf{Client Execution:}
        \Procedure{GradientStep}{$k, w_t$}
        \State $g \gets \nabla f_i(w_t)$ over $\mathcal{P}_k$
        \State \textbf{return} $g$ to server
        \EndProcedure
    \end{algorithmic}
\end{algorithm}

%%%%%%%%%%%%%%%%%%%%%%%%%%%%%%%%%%%%%%%%%%%%%%%%%%%%%%%%%%%%%%%%%%%%%%%%%%%%%%%%%%%%%%%%%%%%%%%%%%%%%%%%%%%%%%%%%%%%%%%%%%%%%%%%%%%%%%%%%%%%%%%%%%%%%%%%%%%%%%
\subsection{Federated Averaging (FedAvg)}
While FedSGD provides a straightforward extension of SGD to the federated setting, it requires frequent communication between clients
and the server, which can be a significant bottleneck in environments with limited bandwidth. To address this issue, the Federated
Averaging (FedAvg) algorithm was proposed, introducing several key innovations to reduce communication overhead while maintaining
model performance.

FedAvg is a generalization of the FedSGD approach by allowing clients to perform multiple local updates using their own data before communicating
with the server. Essentially, this means that each client performs a local gradient descent step upon the his model using local data.
After a number of local gradient updates then clients communicate their model updates instead of the single gradient that is, subsequently
aggregated by the server. Therefore, not only the frequency of communication rounds is reduced, but also local computations are leveraged
more effectively, especially if we take into account the forever increasing performance of contemporary IoT devices utilized as clients.

\begin{algorithm}[H]
    \caption{FederatedAveraging}
    \label{alg:fedavg}
    \begin{algorithmic}[1]
        \State \textbf{Server executes:}
        \State Initialize global model weights $w_0$, and learning rate $\eta$
        \For{each round $t = 1, 2, \ldots$ \textbf{do}}
        \State $m \gets \max(C \cdot K, 1)$
        \State $S_t \gets$ random set of $m$ clients
        \State Send global model $w_t$ to $S_t$ clients
        \For{each client $k \in S_t$ \textbf{in parallel} \textbf{do}}
        \State $w_t^k \gets$ \textsc{ClientUpdate}($k, w_t$)
        \EndFor
        \State $w_{t+1} \gets \sum_{k=1}^K \frac{n_k}{n} w_t^k$
        \EndFor

        \vspace{1em}
        \State \textbf{ClientUpdate}($k, w_t$):
        \State $B \gets$ split $\mathcal{P}_k$ into batches of size $B$
        \For{each local epoch $i$ from 1 to $E$ \textbf{do}}
        \For{batch $b \in B$ \textbf{do}}
        \State $w \gets w - \eta \nabla \ell(w; b)$
        \EndFor
        \EndFor
        \State \textbf{return} $w$ to server
    \end{algorithmic}
\end{algorithm}

All in all, the amount of computation can be adjusted using three key learning parameters: \emph{$C$}, the fraction of clients participating
in each round; \emph{$E$}, the number of local epochs or training passes each client performs over its local dataset per round; and \emph{$B$},
the local minibatch size used for client updates. In the extreme case where $B \rightarrow \infty$ and $E = 1$, FedAvg naturally simplifies to
FedSGD, thereby concealing the unique novelties of FedAvg.
% \begin{enumerate}
% \item \emph{$C$:} The fraction of clients participating in that round.
% \item \emph{$E$:} The number of local epochs-training passes each client makes over its local dataset each round.
% \item \emph{$B$:} The local minibatch size used for client updates
% \end{enumerate


The novelties introduced by FedAvg, briefly outlined above, made the algorithm a more practical and efficient approach to federated learning,
addressing some of the key limitations of FedSGD.\ To sum up, these innovations include reduced communication frequency, local computation
utilization, and improved scalability. By performing multiple local updates before sending the model to the server, FedAvg significantly
reduces the number of communication rounds required. Concurrently, clients leverage their local computational resources more effectively,
due to the increasing processing power of modern edge devices. Finally, by taking load out of the server with the local updates, this
provides the foundations for enhanced scalability as the processing load is distributed.

%%%%%%%%%%%%%%%%%%%%%%%%%%%%%%%%%%%%%%%%%%%%%%%%%%%%%%%%%%%%%%%%%%%%%%%%%%%%%%%%%%%%%%%%%%%%%%%%%%%%%%%%%%%%%%%%%%%%%%%%%%%%%%%%%%%%%%%%%%%%%%%%%%%%%%%%%%%%%%
\subsection{Contempory Challenges of Federated Learning}
Federated Learning (FL) faces several contemporary challenges that need to be addressed to enhance its implementation and efficiency.
The objective of this research was, among else, to provide a pathway on how to overcome some of these issues, yet the entirety of them
would be impossible to be addressed solely with one thesis dissertation. In this spirit,
to be able to illustrate succintly the various shortcomings of Federated Learning, whilst highlighting their relevance to this
research study, an outline of the principal challenges, along with insights into potential solutions that are
presented in our implementation will be presented. The following analysis is based on the journal article introduced by
Li et al.\cite{liFederatedLearningChallenges2020} (2020):
%%%%%%%%%%%%%%%%%%%%%%%%%%%%%%%%%%%%%%%%%%%%%%%%%%%%%%%%%%%%%%%%%%%%%%%%%%%%%%%
\paragraph{Complex infarastracture and Management:} Federated Learning demands a robust and complex infrastructure for coordinating numerous
distributed devices with diverse hardware, software, and network conditions, making setup and maintenance challenging. Ensuring seamless synchronization,
managing device failures, and handling system heterogeneity require sophisticated solutions. Additionally, secure data transmission and efficient
resource allocation are essential. Despite advances in machine learning infrastructure and research, the complexity of FL infrastructure reveals a need
for simplified deployment and management interfaces. This gap highlights the necessity for easy federation of machine learning processes, making FL more
accessible and manageable.
%%%%%%%%%%%%%%%%%%%%%%%%%%%%%%%%%%%%%%%%%%%%%%%%%%%%%%%%%%%%%%%%%%%%%%%%%%%%%%%
\paragraph{Expensive Communication:} Communicaion Overhead is a critical bottleneck in Federated Networks, that necessitate against
transmitting locally generated data. The reality in Federated Learning Setups differs significantly from the almost perfect communication
setups of traditional Centralized Learnig Data centers, where stable and incredibly fast connections are assumed with restricted faillures.
In contrast, Federated Learning is conducted on a direct opposite manner, meaning that learning is based on utilizing a vast number of
devices, like smartphones, resulting in substantially reduced communication capabilities with several limiting factors including bandwidth
energy and transmission power. Therefore, in edge FL, where the datasets are small and the connections are weak and unreliable the communication
to computation ration skyrockets.
%%%%%%%%%%%%%%%%%%%%%%%%%%%%%%%%%%%%%%%%%%%%%%%%%%%%%%%%%%%%%%%%%%%%%%%%%%%%%%%
\paragraph{System Heterogenity:} Devices in federated networks exhibit significant variability in storage, computational power, and communication
capabilities due to differences in hardware (CPU, memory), network connectivity (3G, 4G, 5G, Wi-Fi), and power levels (battery). Typically, only
a small fraction of devices are active at any given time, and devices may drop out due to connectivity or energy constraints. These system-level
disparities aggravate challenges such as managing stragglers, who delay updates, and ensuring fault tolerance within the network. Robust methods
are needed to handle low participation and heterogeneous environments effectively.
%%%%%%%%%%%%%%%%%%%%%%%%%%%%%%%%%%%%%%%%%%%%%%%%%%%%%%%%%%%%%%%%%%%%%%%%%%%%%%%
\paragraph{Statistical Data Heterogeneity:} The need for independent and identically distributed (i.i.d.) datasets is crucial for efficiency in Distributed
Systems. In Federated Systems, while non-i.i.d.\ data can be used, it significantly impacts the speed and difficulty of convergence. Each client trains
on its local dataset, adding complexity and likely resulting in skewed updates based on user preferences. The server, due to privacy concerns, cannot
know the dataset distribution, necessitating sophisticated approaches to maintain accuracy while learning from highly biased datasets. These methods
must effectively address the challenges posed by data heterogeneity to ensure efficient federated learning.
%%%%%%%%%%%%%%%%%%%%%%%%%%%%%%%%%%%%%%%%%%%%%%%%%%%%%%%%%%%%%%%%%%%%%%%%%%%%%%%
\paragraph{Privacy and Security:} The main objective of Federated Learning is to ensure client privacy. However, shared parameters can be exploited by
malicious participants or third parties to extract sensitive information, undermining this goal. It is commonly assumed that all clients act in good faith,
but some may be malicious, attempting to compromise the model's integrity. They can perform model poisoning or data poisoning attacks, using corrupted data
to degrade model accuracy or introduce backdoors. Therefore, robust security measures are crucial to maintain the integrity and privacy of the FL process.

%%%%%%%%%%%%%%%%%%%%%%%%%%%%%%%%%%%%%%%%%%%%%%%%%%%%%%%%%%%%%%%%%%%%%%%%%%%%%%%%%%%%%%%%%%%%%%%%%%%%%%%%%%%%%%%%%%%%%%%%%%%%%%%%%%%%%%%%%%%%%%%%%%%%%%%%%%%%%%
\subsection{Countermeasures Introduced in Thesis}
This thesis addresses several key challenges in Federated Learning, some as primary objectives and others as secondary considerations.
The following discussion is organized by relevance to the thesis contents, starting with the main objectives: mitigating communication overhead and providing a
user-friendly interface for FL experimentation.

To reduce communication overhead in FL, this work employs strategies such as minimizing the size of transmitted messages, reducing
the number of communication rounds, and limiting communication frequency. Specifically, the thesis implements the Functional Dynamic
Averaging (FDA) synchronization technique (\cref{sec:synch_techniques}) to enhance convergence speed, effectively balancing the trade-offs
between communication and computation. Additionally, the Flower Framework (\cref{sec:flower}) is utilized to facilitate a seamless transition
from traditional ML frameworks to FL setups, offering an accessible and efficient infrastructure.

While the issue of system heterogeneity is not directly addressed, the thesis integrates a Federated Learning monitoring framework with the Ns3
network simulator (\cref{sec:ns3}). This integration allows for experimentation with various communication setups and testing different device
models with diverse characteristics (e.g., power consumption, computational capabilities). Monitoring data heterogeneity is also part of the
framework to assess system performance under different conditions.

Privacy and security concerns, although not directly tackled in this study, are acknowledged as critical areas for future research. The experiments
are conducted in a controlled environment, which minimizes the immediate need for additional privacy measures. However, addressing these concerns
remains essential for the advancement of Federated Learning.

%%%%%%%%%%%%%%%%%%%%%%%%%%%%%%%%%%%%%%%%%%%%%%%%%%%%%%%%%%%%%%%%%%%%%%%%%%%%%%%
%%%%%%%%%%%%%%%%%%%%%%%%%%%%%%%%%%%%%%%%%%%%%%%%%%%%%%%%%%%%%%%%%%%%%%%%%%%%%%%%%%%%%%%%%%%%%%%%%%%%%%%%%%%%%%%%%%%%%%%%%%%%%%%%%%%%%%%%%%%%%%%%%%%%%%%%%%%%%%
